Trong suốt quá trình thực hiện đề tài tốt nghiệp phát triển trò chơi hành động phiêu lưu 2D "Dream Witch", Em đã nghiên cứu, thiết kế và cài đặt một số giải pháp kỹ thuật cốt lõi nhằm giải quyết các bài toán khó về mặt hiệu năng vận hành, quản lý logic máy trạng thái nhân vật phức tạp, cũng như kiểm soát luồng hội thoại phi tuyến của cốt truyện. Chương này sẽ trình bày chi tiết ba giải pháp đóng góp nổi bật và tâm đắc nhất của đề tài.

\section{Tối ưu hóa hiệu năng bằng kỹ thuật Distance Culling tùy biến}

\subsection{Đặt vấn đề và bài toán tối ưu}
Trong thể loại game đi cảnh thế giới mở đa tuyến (Metroidvania), môi trường trò chơi thường là một lưới phòng bản đồ liền mạch chứa hàng ngàn vật thể động như quái vật thường (Enemy), Boss lớn, bẫy môi trường di động, các animator và diễn họa chuyển động dạng xương (Spine 2D skeletal animation) đối với các thực thể Boss lớn. Với cấu trúc truyền thống của Unity, mọi thực thể được kích hoạt trên Scene đều sẽ liên tục thực thi các hàm cập nhật logic định kỳ như \texttt{Update()} hoặc \texttt{FixedUpdate()} để tính toán vật lý, phát hiện va chạm và chạy các giải thuật AI định tuyến (A* Pathfinding). 

Đặc biệt, đối với các Boss khổng lồ, Spine 2D skeletal animation đòi hỏi chi phí tính toán CPU rất lớn cho việc thực hiện phép toán biến dạng mạng lưới đỉnh (Mesh Deformation) và nội suy khung xương theo thời gian thực. Đối với nhân vật chính và quái vật thường, việc chạy đồng thời hệ thống Animator với số lượng lớn trạng thái chuyển tiếp cũng gây áp lực không nhỏ lên CPU. Việc vận hành đồng thời toàn bộ các tài nguyên này sẽ dẫn đến hiện tượng quá tải CPU (CPU Bottleneck), làm sụt giảm tốc độ khung hình (khung hình giảm xuống dưới 30 FPS hoặc gây trễ đầu vào) trên các máy tính có cấu hình thấp hoặc trung bình. Do đó, bài toán đặt ra là cần xây dựng một giải pháp kỹ thuật để cô lập và vô hiệu hóa một cách chọn lọc các thực thể nằm ngoài tầm nhìn của người chơi mà không làm ảnh hưởng đến tính nhất quán của trạng thái thế giới hoặc gây cảm giác gián đoạn cho trải nghiệm chơi game.

\subsection{Giải pháp hiện thực hóa}
Đề tài đề xuất giải pháp thiết kế một hệ thống quản lý culling khoảng cách tùy biến dựa trên hai lớp cốt lõi là \texttt{CullingManager} và \texttt{DistanceCullingTarget} (như đã giới thiệu sơ bộ ở Chương 5). Thay vì sử dụng hệ thống culling hình nón camera truyền thống (Camera Frustum Culling - vốn chỉ ngắt tiến trình render nhưng vẫn chạy CPU logic cập nhật và tính toán vật lý), giải pháp này thực hiện ngắt chọn lọc cả logic cập nhật và hiển thị dựa trên khoảng cách vật lý thực tế.

Cụ thể, thuật toán kiểm tra khoảng cách được định nghĩa dựa trên khoảng cách Euclid giữa vị trí của người chơi $P(x_p, y_p)$ và vị trí của đối tượng culling $T(x_t, y_t)$:
\begin{equation}
d = \sqrt{(x_p - x_t)^2 + (y_p - y_t)^2}
\end{equation}

Mỗi đối tượng thuộc diện culling (\texttt{DistanceCullingTarget}) được phân loại thành 3 nhóm mục tiêu chính: \textbf{Quái vật (Enemy)}, \textbf{Đạn bay (Projectile)} và \textbf{Hiệu ứng hình ảnh (Vfx)}. Thay vì dùng một bán kính cứng nhắc, hệ thống áp dụng \textbf{cơ chế khoảng trễ Hysteresis} với hai ngưỡng: \textbf{Ngưỡng bật (mặc định 18f)} để kích hoạt lại đối tượng và \textbf{Ngưỡng tắt (mặc định 22f)} để vô hiệu hóa đối tượng, nhằm tránh hiện tượng bật/tắt liên tục khi người chơi đứng ở ranh giới. Các đối tượng đặc thù như Boss lớn có thể ghi đè hai ngưỡng này thành các giá trị độc lập.

Quy trình xử lý culling diễn ra như sau:
\begin{enumerate}
    \item \texttt{CullingManager} phân tách dữ liệu lưu trữ thành hai nhóm để tối ưu truy xuất: sử dụng \texttt{List<DistanceCullingTarget>} để lưu danh sách các đối tượng chịu culling khoảng cách nhằm tăng tốc độ lặp (iteration) theo frame, và sử dụng \texttt{HashSet<RoomController>} để lưu trữ tập hợp các phòng đang hoạt động nhằm đạt độ phức tạp $O(1)$ khi kiểm tra trạng thái phòng.
    \item Định kỳ chạy coroutine \texttt{CullingLoop} sau mỗi $0.2$ giây (chu kỳ kiểm tra cố định để tiết kiệm năng lượng CPU thay vì cập nhật từng frame), \texttt{CullingManager} tiến hành duyệt danh sách và phối hợp culling phân tầng.
    \item Quá trình kiểm tra diễn ra theo trình tự ưu tiên:
    \begin{itemize}
        \item \textbf{Ưu tiên Phòng (Room Sleeping):} Nếu đối tượng thuộc về một phòng ngủ đông (thông qua lớp cầu nối \texttt{RoomCullingMember} báo cờ \texttt{isRoomSleeping = true}), nó lập tức bị ép vào trạng thái tắt mà không cần xét khoảng cách, ngăn chặn tuyệt đối lỗi AI quái vật vô tình thức dậy.
        \item \textbf{Ưu tiên Khoảng cách:} Nếu phòng đang hoạt động, hệ thống mới tiến hành xét khoảng cách Euclid $d$. Nếu $d >$ \texttt{Disable Distance}, đối tượng chuyển sang trạng thái \textit{Culled} (tắt). Ngược lại, nếu $d <$ \texttt{Enable Distance}, đối tượng chuyển sang \textit{Active} (bật).
    \end{itemize}
    \item Khi chuyển sang trạng thái \textit{Culled}, hệ thống ngắt chọn lọc các thành phần tài nguyên nặng:
    \begin{itemize}
        \item Vô hiệu hóa component tính toán vật lý (\texttt{Rigidbody2D}, \texttt{Collider2D}) để giải phóng tài nguyên Engine (Physics2D).
        \item Vô hiệu hóa Animator đối với quái thường hoặc Spine component đối với Boss để dừng tính toán hoạt ảnh.
        \item Dừng các luồng xử lý AI tìm đường A* (\texttt{AIPath} và \texttt{Pathfinder} components).
    \end{itemize}
    \item Khôi phục tức thời các thành phần trên về trạng thái hoạt động ban đầu mà không có độ trễ hiển thị đối với người chơi khi chuyển về trạng thái \textit{Active}.
\end{enumerate}

\subsection{Camera Frustum Culling cho Tilemap}
Bên cạnh việc culling các thực thể động, hệ thống còn mở rộng tối ưu hóa culling tĩnh cho các \texttt{TilemapRenderer} chiếm dung lượng lớn bằng mặt phẳng chóp cụt camera (Camera Frustum Culling).
Trong mỗi \texttt{RoomController}, hệ thống bắt sự kiện camera di chuyển để tự động tính toán mặt phẳng chóp cụt góc nhìn hiện tại qua hàm \texttt{GeometryUtility.}\allowbreak\texttt{CalculateFrustumPlanes}. Tiếp theo, hệ thống dùng phương thức kiểm tra giao cắt AABB (\texttt{GeometryUtility.}\allowbreak\texttt{TestPlanesAABB}) giữa chóp cụt camera với \texttt{Bounds} của Tilemap, cộng thêm khoảng đệm \texttt{cameraVisibilityPadding}. Nếu Tilemap nằm ngoài camera, \texttt{Renderer.enabled} sẽ bị tắt, giúp giảm triệt để lệnh Draw Calls đồ họa tĩnh.

\subsection{Kết quả đạt được}
Giải pháp culling tùy biến đã mang lại kết quả tối ưu hóa hiệu năng vượt trội. \textbf{Tốc độ khung hình (FPS) tăng trưởng ổn định:} Khi thám hiểm các phân khu có mật độ quái vật dày đặc, FPS trung bình tăng từ 42 lên mức 60 ổn định, loại bỏ hoàn toàn giật lag cục bộ. 

\textbf{Giảm tải CPU đáng kể:} Tổng thời gian CPU xử lý cập nhật hoạt ảnh giảm 75\% (từ 12ms xuống dưới 3ms mỗi frame). Lực cản vật lý và va chạm thừa của đối tượng ngoài tầm nhìn được triệt tiêu, giảm tải overhead cho Physics Engine tới 80\%. Đồng thời, nhờ \textbf{hiệu ứng hồi phục trơn tru} với buffer khoảng cách thông minh, người chơi không cảm nhận được hiện tượng bật/tắt (pop-in), đảm bảo trải nghiệm khám phá liền mạch tuyệt đối.

\section{Kiến trúc máy trạng thái phân tầng cho hệ thống nhân vật phức tạp}

\subsection{Đặt vấn đề và bài toán thiết kế}
Nhân vật chính trong game đi cảnh hành động đòi hỏi tập hợp các hành vi di chuyển cực kỳ đa dạng và phản hồi nhịp độ cao bao gồm: chạy trên mặt đất, nhảy đơn, nhảy đôi, nhảy bám tường, lướt đa hướng (Dash), leo vách đá và bám tường (Wall Climb), kết hợp với chuỗi combo tấn công cận chiến và cơ chế né lướt, chặn đỡ đòn (Dodge/Guard). Đồng thời, nhân vật có khả năng biến hình (Form Switch) sang các hình thái khác nhau dẫn đến việc thay đổi thuộc tính vật lý (như trọng lực, lực nhảy, tốc độ di chuyển) và tập hợp kỹ năng thi triển.

Nếu cài đặt tất cả các logic này trong một lớp duy nhất sử dụng các câu lệnh điều kiện lồng nhau (\texttt{if-else} hoặc \texttt{switch-case} lớn), mã nguồn sẽ nhanh chóng trở nên cực kỳ phức tạp (spaghetti code), rất khó bảo trì và phát sinh hàng loạt lỗi logic nghiêm trọng về trạng thái điều khiển (ví dụ: nhân vật vẫn có thể thực hiện combo tấn công khi đang leo tường hoặc đang lướt tự do). Vì vậy, bài toán đặt ra là cần thiết kế một kiến trúc điều khiển tách biệt, an toàn về trạng thái và có khả năng mở rộng tốt.

\subsection{Giải pháp hiện thực hóa}
Đề tài đã thiết kế và cài đặt kiến trúc máy trạng thái phân tầng (\textit{Hierarchical State Machine - HSM}) kết hợp mô hình hướng đối tượng chặt chẽ. Hệ thống được xây dựng xung quanh ba lớp chính:
\begin{itemize}
    \item \texttt{PlayerController}: Lớp điều phối trung tâm giữ vai trò trung gian chứa các tham chiếu tài nguyên vật lý, bộ quản lý hoạt ảnh của nhân vật, thông số thuộc tính và đối tượng máy trạng thái.
    \item \texttt{PlayerStateMachine}: Quản lý trạng thái hiện tại và điều phối quá trình chuyển dịch trạng thái thông qua phương thức \texttt{Initialize()} và \texttt{Change\-State()}.
    \item \texttt{PlayerState}: Lớp cơ sở trừu tượng cho tất cả các trạng thái, cung cấp vòng đời chuẩn hóa: \texttt{Enter()} (khởi chạy trạng thái), \texttt{LogicUpdate()} (cập nhật logic), \texttt{PhysicsUpdate()} (cập nhật vật lý) và \texttt{Exit()} (dọn dẹp).
\end{itemize}

Các trạng thái cụ thể như \texttt{IdleState}, \texttt{MoveState}, \texttt{JumpState}, \texttt{DashState}, và \texttt{WallClimbState} được kế thừa từ \texttt{PlayerState}. Để giải quyết triệt để bài toán biến đổi hình dạng (Form Switch) mà không làm gián đoạn máy trạng thái, giải pháp tích hợp thêm phân hệ \texttt{PlayerFormManager} hoạt động song song.

Quy trình xử lý Form Switch diễn ra như sau:
\begin{enumerate}
    \item Khi người chơi nhấn nút chuyển dạng, \texttt{PlayerFormManager} truy vấn dữ liệu từ tệp tài nguyên cấu hình \texttt{PlayerFormDataSO} (ScriptableObject) của Form đích.
    \item Thay vì phá hủy máy trạng thái và tạo mới, hệ thống thực hiện cập nhật lại tập hợp các tham chiếu hoạt ảnh Unity Animator (thông qua cơ chế nạp chuỗi ký tự string-based trực tiếp) và các chỉ số thuộc tính runtime hiện hành (Max HP, tốc độ chạy, lực nhảy) ngay bên trong đối tượng \texttt{PlayerController}.
    \item Trạng thái hiện tại trong máy trạng thái (\texttt{currentState}) nhận diện sự thay đổi thuộc tính và tự động điều chỉnh lực vật lý áp dụng trong hàm \texttt{PhysicsUpdate()} tiếp theo, duy trì trạng thái di chuyển hiện tại một cách liên tục và mượt mà.
\end{enumerate}

\subsection{Kết quả đạt được}
Kiến trúc máy trạng thái phân tầng mang lại những cải tiến đáng kể. \textbf{Triệt tiêu lỗi xung đột logic:} Đóng gói logic vào từng lớp giúp kiểm soát tuyệt đối luồng nhập liệu, ngăn chặn các hành vi bị cấm (ví dụ: cấm tấn công khi đang leo tường do lớp \texttt{WallClimbState} không chứa lệnh tấn công).

\textbf{Độ linh hoạt và mở rộng cao:} Việc thêm trạng thái hay hình thái biến hình mới chỉ cần tạo lớp kế thừa \texttt{PlayerState} hoặc tệp \texttt{PlayerFormDataSO} mới mà không phải sửa mã nguồn cũ, tuân thủ chặt chẽ nguyên tắc đóng/mở (Open/Closed). Nhờ đó, \textbf{hiệu ứng chuyển đổi diễn ra mượt mà}, gán Animator Controller mới tức thì với độ trễ bằng 0, mang lại cảm giác chiến đấu nhịp độ cao sống động.

\section{Giải pháp tích hợp Yarn Spinner và quản lý máy trạng thái thế giới phi tuyến}

\subsection{Đặt vấn đề}
Trò chơi phiêu lưu nhập vai 2D đòi hỏi một hệ thống hội thoại phong phú, đa tuyến và phụ thuộc chặt chẽ vào tiến trình thế giới của người chơi (Quest progression). Ví dụ, NPC chủ cửa hàng ở làng (Village) sẽ thay đổi nội dung hội thoại, mở khóa các mặt hàng mới hoặc giao nhiệm vụ mới dựa trên việc người chơi đã thu thập được các cuốn sách phép (\texttt{SpellBooks}) nào, đã kích hoạt cổng dịch chuyển (\texttt{Portals}) nào hoặc đã tiêu diệt các Boss cụ thể nào.

Việc tự phát triển một hệ thống hội thoại có khả năng rẽ nhánh phức tạp từ đầu đòi hỏi khối lượng công việc lập trình giao diện lớn và dễ phát sinh lỗi khi hiển thị đa ngôn ngữ (Localization) hoặc các ký tự đặc biệt. Đồng thời, việc liên kết luồng hội thoại với trạng thái thế giới động để lưu trữ vĩnh viễn tiến trình của người chơi vào tệp JSON là một bài toán khó đòi hỏi thiết kế kiến trúc lưu trữ nhất quán.

\subsection{Cơ chế tích hợp và đồng bộ hóa trạng thái}
Đề tài đã giải quyết bài toán này bằng cách tích hợp trực tiếp thư viện Yarn Spinner cho Unity kết hợp với giải pháp tuần tự hóa dữ liệu đồng bộ trạng thái lưu trữ (\texttt{GameSaveManager}).

Giải pháp kỹ thuật bao gồm các thành phần:
\begin{itemize}
    \item \textbf{Kịch bản hội thoại Yarn Script:} Các cuộc hội thoại được viết trên các tệp \texttt{.yarn} độc lập bằng cú pháp trực quan. Các điều kiện rẽ nhánh được biểu diễn bằng cấu trúc biến số nội tại của Yarn (ví dụ: \texttt{\$has\_defeated\_golem} hoặc \texttt{\$has\_unlocked\_fire\_spell}).
    \item \textbf{Lớp cầu nối YarnDialogueView:} Lớp này chịu trách nhiệm hiển thị giao diện UI hội thoại dạng văn bản hỗ trợ nhiều ngôn ngữ và các lựa chọn rẽ nhánh cho người chơi.
    \item \textbf{Cơ chế đồng bộ hóa biến số (Variable Storage Synchronizer):} Khi trò chơi bắt đầu hoặc tải lại tiến trình, lớp \texttt{GameSaveManager} sẽ đọc dữ liệu từ tệp lưu trữ JSON (\texttt{GameSaveData}) chứa thông tin tiến trình thực tế của người chơi (bao gồm các thuộc tính như các Boss đã hạ gục \texttt{bossDefeats}, các cửa đã mở \texttt{doors}, cổng dịch chuyển đã đi qua \texttt{portals}, và tiến trình cây kỹ năng \texttt{skillTreeProgress}). Hệ thống sau đó tự động chuyển đổi các trạng thái này và nạp vào bộ nhớ biến số của Yarn Spinner thông qua phương thức \texttt{VariableStorage.\allowbreak SetValue()}.
\end{itemize}

Ngược lại, khi người chơi thực hiện các lựa chọn hội thoại làm thay đổi biến tiến trình trong Yarn Spinner, các sự kiện callback từ Yarn Spinner sẽ lập tức gọi các hàm thay đổi trạng thái tương ứng trong \texttt{GameSaveManager} để cập nhật lại tệp tin lưu trữ JSON thời gian thực, đảm bảo tiến trình hội thoại và tiến trình thế giới luôn đồng bộ tuyệt đối.

\subsection{Kết quả đạt được}
Giải pháp tích hợp và quản lý trạng thái hội thoại phi tuyến đã mang lại những kết quả tích cực:
\begin{itemize}
    \item \textbf{Quy trình viết kịch bản chuyên nghiệp:} Tách biệt hoàn toàn công việc viết kịch bản hội thoại ra khỏi mã nguồn C\#. Nhà thiết kế game có thể cập nhật nội dung lời thoại với nhiều ngôn ngữ trực tiếp trên các tệp \texttt{.yarn} mà không cần biên dịch lại dự án Unity.
    \item \textbf{Lưu giữ tiến trình đồng bộ:} Trạng thái hội thoại và tiến trình thế giới được lưu trữ đồng bộ vào tệp JSON. Khi tải lại game, NPC hiển thị đúng lời thoại tương ứng với trạng thái nhiệm vụ đã lưu, loại bỏ lỗi reset tiến trình.
    \item \textbf{Hỗ trợ hội thoại nhiều ngôn ngữ khác nhau:} Hệ thống hoạt động ổn định, có khả năng tích hợp và hiển thị văn bản đa ngôn ngữ (Localization) mượt mà, kết hợp rẽ nhánh trực quan tương thích tốt với điều khiển đầu vào.
\end{itemize}